% TOP_PROBLEM: {"id": 3078, "title": "Simplexity of the n-cube", "category_id": 6, "category": {"id": 6, "name": "geometry", "display_name": "Geometry", "description": "Euclidean and non-Euclidean geometry, geometric structures.", "slug": "geometry", "order_index": 6, "created_at": "2026-07-31T15:26:25.671Z"}, "status": "open", "problem_number": "OPG-1820", "set_id": 12, "statusQualification": "T(n) permits additional vertices and non-face-to-face intersections; the stricter vertex triangulation number T*(n) is distinguished explicitly."}
% FIRST_POSTED: 2026-10-01
% ATTEMPT_STATUS: unresolved
% UnsolvedMath ID 3078; source catalogue code OPG-1820
% !TeX program = lualatex
% GPT-6 Astra Ultra research attempt; independent partial work, 2026-10-01.
\documentclass[11pt,a4paper]{article}
\usepackage[margin=25mm]{geometry}
\usepackage{fontspec}
\setmainfont{lmroman10-regular.otf}[BoldFont=lmroman10-bold.otf,
 ItalicFont=lmroman10-italic.otf,BoldItalicFont=lmroman10-bolditalic.otf]
\usepackage{amsmath,amssymb,mathtools}
\usepackage{unicode-math}
\setmathfont{latinmodern-math.otf}
\AtBeginDocument{\let\setminus\smallsetminus}
\usepackage{xurl,hyperref}
\hypersetup{colorlinks=true,urlcolor=blue}
\setcounter{secnumdepth}{0}
\setlength{\parindent}{0pt}
\setlength{\parskip}{5pt}
\setlength{\emergencystretch}{3em}
\begin{document}
\section*{Simplexity of the n-cube}\label{3078}
{\small UnsolvedMath ID 3078 · OPG-1820 · Geometry · OpenGarden}
\subsection{Definitions and mathematical statement}
Let $n\ge1$ and let $Q_n=[0,1]^n\subset\mathbb R^n$.
An $n$-simplex is the convex hull of $n+1$ affinely independent points.
A decomposition of $Q_n$ means a finite family
$\mathcal D=\{\Delta_1,\ldots,\Delta_m\}$ of closed $n$-simplices
contained in $Q_n$, with union $Q_n$ and pairwise disjoint interiors.
Define
\[
 T(n)=\min\{m:Q_n\text{ has such a decomposition into }m\text{ simplices}\}.
\]
Determine $T(n)$ as a function of dimension, by matching constructive
upper bounds with lower bounds valid for every allowed decomposition.

The convention here follows the source's explicit definition of $T$:
vertices of the simplices may be arbitrary points of the cube, and
intersections of distinct simplices need not be entire common faces.
For comparison, its separate parameter $T^*(n)$ imposes both additional
requirements that every simplex vertex be a cube vertex and that the
simplices form a face-to-face triangulation. Thus $T(n)\le T^*(n)$, and
these quantities must not be silently identified.

For any permutation $\sigma$ of $\{1,\ldots,n\}$, the region
$0\le x_{\sigma(1)}\le\cdots\le x_{\sigma(n)}\le1$ is an $n$-simplex.
The resulting $n!$ simplices give a decomposition, ensuring the minimization
is over a nonempty class. Scaling or translating the cube does not
alter either minimum. The target counts full-dimensional pieces; lower-dimensional
faces and their subdivisions are not counted as additional simplices.
\subsection{Short English statement}
What is the minimum number of full-dimensional simplices needed to decompose a cube, under the source’s unrestricted decomposition convention?
\subsection{Research attempt: a volume bound and unrestricted three-dimensional optimality}
\paragraph{Scope and literature check.}
GPT-6 Astra Ultra, 1 October 2026. Additional vertices and intersections
which are not common whole faces are allowed throughout the lower
bounds. The original discussion [S2] distinguishes $T$ from $T^*$,
reports the familiar five-tetrahedron cube decomposition, and points to
stronger asymptotic constructions and lower bounds. We reconstruct
elementary certificates here; none is claimed as a new record.

\paragraph{1. A volume lower bound for arbitrary simplex vertices.}
For vertices $v_0,\ldots,v_n\in[0,1]^n$, form the matrix with rows
$(1,v_i)$. Its determinant has absolute value $n!$ times the simplex
volume. Subtract the first column from twice each of the other columns.
The resulting rows $(1,2v_i-\boldsymbol1)$ have Euclidean norms at
most $\sqrt{n+1}$, and the determinant is multiplied by $2^n$.
The determinant bound by the product of row norms follows from
Gram--Schmidt orthogonalization: orthogonal projection cannot increase
any successive row length. Hence every allowed simplex has volume at most
\[
 V_n\le\frac{(n+1)^{(n+1)/2}}{2^n n!},\qquad
 T(n)\ge\left\lceil\frac{2^n n!}{(n+1)^{(n+1)/2}}\right\rceil.
\]
Volumes add because interiors are disjoint and boundaries have
$n$-dimensional measure zero. The argument does not assume a
face-to-face triangulation or vertices at cube corners.

\paragraph{2. Construct five tetrahedra in the three-cube.}
The four even-parity corners $000,110,101,011$ form a central
tetrahedron of determinant magnitude two and volume $1/3$. At each
odd-parity corner $100,010,001,111$, take that corner together with
its three adjacent even-parity corners. These are four tetrahedra of
volume $1/6$ each. Their interiors correspond respectively to
\[
 x>y+z,\quad y>x+z,\quad z>x+y,\quad x+y+z>2
\]
inside the cube. No two inequalities can hold together: two among
the first three contradict nonnegativity, and the first together
with the last implies $2x>2$, with analogous contradictions for
$y,z$. The complement of these four corner interiors is the central
tetrahedron, described by the reversed weak inequalities. Alternatively,
its barycentric coordinates are
$1-(x+y+z)/2$, $(x+y-z)/2$, $(x+z-y)/2$, $(y+z-x)/2$.
They sum to one and are all nonnegative exactly there. This proves
coverage and disjoint interiors, so $T(3)\le5$.

\paragraph{3. Four tetrahedra are impossible even with extra vertices.}
Every square boundary face of the cube must be covered by triangular
faces of the tetrahedra. A supporting cube plane intersects a
full-dimensional tetrahedron in a face, so an intersection of positive
area is an entire triangular face, not an arbitrary polygon. At least
two such triangles are needed for each square: a triangle contained
in a square cannot contain all four extreme corners. Thus there are
at least twelve tetrahedron faces lying on the cube boundary.

A tetrahedron has at most three such boundary faces. Two faces of
a tetrahedron meet in an edge, so cannot lie on two opposite parallel
cube planes. There are only three pairs of coordinate planes, allowing
at most one face per pair. If at most four tetrahedra covered the
cube, the count forces exactly four tetrahedra with three boundary
faces each. The three supporting coordinate planes for any such
tetrahedron meet at a cube corner; its remaining three vertices lie
on the three incident cube edges. Their distances from that corner
are at most one, so its volume is at most $1/6$. The total volume
would be at most $4/6<1$, a contradiction. Therefore $T(3)=5$.

\paragraph{Verification and remaining gap.}
The boundary count tolerates subdivisions and nonmatching interior
faces. Lower-dimensional intersections with cube faces do not cover
positive-area portions and do not count as triangles. The determinant
checks give volumes $1/3$ and four copies of $1/6$, summing to one.
In higher dimension the universal determinant estimate is generally
far from the constructive $n!$ bound in the definition. Large-volume
simplices cannot simply be packed without overlap to attain that
estimate. Controlling this packing compatibility is the missing step.

\paragraph{Reproduction of finite checks.}
The finite checks stated above were run with Python 3 using the repository
script [C1]. They verify the stated examples and finite ranges only.

\subsection{Final status}
UNRESOLVED as a function of arbitrary dimension. The unrestricted
volume lower bound and a complete proof of the known exact value
$T(3)=5$ are supplied. Equality of $T(n)$ and $T^*(n)$ is not assumed,
and no novel optimal construction or global completion claim is made.

\subsection{Sources}
\begin{itemize}
\item[S1] ULAM.AI and the UnsolvedMath Contributors, supplied problems.json, record 3078 (OPG-1820), OpenGarden collection. Catalogue identity and source transcription; CC BY 4.0.
\url{https://huggingface.co/datasets/ulamai/UnsolvedMath}.
\item[S2] Open Problem Garden, Simplexity of the n-cube. Original problem and discussion; the mathematical formulation below is expanded from the supplied source record.
\url{http://www.openproblemgarden.org/op/simplexity_of_the_cube}.
\item[S3] Open Problem Garden, \emph{Simplexity of the n-cube}, original maintained problem discussion including the distinction between $T$ and $T^*$ and its Haiman/Smith bibliography. HTTPS version checked 1 October 2026.
\url{https://www.openproblemgarden.org/op/simplexity_of_the_cube}.
\item[C1] Reproducible finite-check code, executed 1 October 2026 with Python 3 (standard library only):
\path{scripts/check_attempt_batch_geometry_2026_10_01.py}.
\end{itemize}
\end{document}
